\section{Assumption 3}\label{s:ass3}

\lem{4}{For every agent capable of intentional action, there is at least one project of objective worth in which she is able to engage.}

\paragraph{The sentence as a whole.} Assumption~3 carries the weight of the availability claim. Its logical form is $\forall a\,\exists p\,(W(p) \wedge E(a,p))$: for every agent there is a project that is of objective worth and in which she is able to engage. Each conjunct of the matrix makes its own demand. The first requires that the project have worth in the restricted sense of Assumption~2. The second requires that the agent be able to engage in it. Neither alone would support Theorem~4: that there are projects of objective worth says nothing about any particular agent, and that every agent can engage in some project says nothing about worth, since the project might be the counting of grass.

EXI-01-c calls this the premise a critic will press, and it is. The comment examines its quantifiers, its domain, the modality of \enquote{able}, and the support it receives in the family.

\lem{4.1}{For every agent}

\paragraph{The universal quantifier.} The quantifier is unrestricted within its domain: it covers every agent capable of intentional action, without exception for circumstance, talent, wealth, health or location. This is what makes the conclusion of the family universal, and it is what makes Assumption~3 vulnerable, since one counterexample suffices to refute it.

\paragraph{The order of the quantifiers.} The universal quantifier precedes the existential one.

\begin{proposition}[Order of quantifiers]\label{p:order}
Assumption~3 does not imply that there is a single project of objective worth in which every agent is able to engage, and Theorem~4 does not require it.
\end{proposition}

\begin{proof}
The form $\forall a\,\exists p\,\phi(a,p)$ does not imply $\exists p\,\forall a\,\phi(a,p)$: each agent may have her own project. The proof of Theorem~4 instantiates the existential quantifier for each agent separately (\Lm{6}), and so uses only the weaker form.
\end{proof}

The proposition marks the difference between EXI-01-a and views on which meaning is available to all because there is one project open to all. On some religious views, the love of God is such a project; on some moral views, the moral life is. EXI-01-a does not need any such project. The prisoner's project and the scholar's may have nothing in common except that each has objective worth and each is within the reach of the one who engages in it. This is one reason why the availability claim is more plausible than it might seem: it does not require that the same opportunity be open to everyone, only that each person have some opportunity.

The stronger form is nonetheless arguably true, and EXI-01-c comes close to asserting it: its Assumption~2.2 says that every agent can direct caring activity at another person or at some part of the world she can come to understand. That is a disjunction of two kinds of project rather than a single project, but the kinds are so general that nearly every agent has both. The text does not need this, and the Commentary does not assert it.

\lem{4.2}{capable of intentional action}

\paragraph{The domain.} The domain of the quantifier is agents capable of intentional action. The phrase has a standard sense in the philosophy of action: an intentional action is one done for a reason, one to which the question \enquote{why?}, in the sense that asks for the agent's reason, applies \citep{Anscombe1957}. An agent capable of intentional action is one who can act for reasons. The restriction is natural, since engagement, on the conative reading of \Lm{2.5}, involves acting in ways directed at an object out of care for it, which is acting for a reason.

The same phrase defines the domain of Lemma~3.1 of EXI-04-a, which holds that any agent capable of intentional action takes some ends seriously, and of the theorems of EXI-01-c and EXI-01-d. The agreement is not accidental. The discipline treats the capacity for intentional action as the threshold at which questions of meaning and absurdity arise: below it there is no seriousness to collide with detachment (EXI-04-a) and no engagement to constitute meaning (EXI-01-a).

\paragraph{Capacity and exercise.} The phrase speaks of capacity, not exercise. An agent who is asleep, or unconscious under anesthesia, remains capable of intentional action and remains in the domain. An agent in a persistent vegetative state, if she has permanently lost the capacity, is outside it. The boundary is not sharp, and the text does not claim that it is: capacities for intentional action come in degrees and develop over time, as in children, and decline, as in dementia. The availability claim holds for each agent while she has the capacity, and EXI-01-c notes, in reply to its third objection, that an agent for whom nothing worthwhile is within reach of attention is, if possible at all, a case of the absence of agency rather than of its exercise.

\paragraph{\enquote{Everyone}.} The family conclusion says that meaning is available to \enquote{everyone}, and Remark~5 says that the theorem answers the question \enquote{whether meaning is available to everyone} (\Lm{9.2}). The theorem itself says \enquote{every such agent}. The two are reconciled by reading \enquote{everyone} as \enquote{every agent capable of intentional action}, and the reading is the natural one: the question whether meaning is available to someone is a question about someone who could engage with something. Human beings who are not, or not yet, or no longer, agents in this sense raise a different question, whether their lives can have meaning in a sense other than the one Definition~1 defines, for instance through the engagement of others with them. Volume~III, Part~A, takes up that question; the family conclusion does not answer it and does not need to.

\lem{4.3}{there is at least one project of objective worth}

\paragraph{The minimal claim.} \enquote{At least one} states the minimum that Theorem~4 needs. A single project of objective worth in which an agent can engage suffices for a life of positive meaning to be available to her. The assumption says nothing about how many such projects there are, how great their worth is or how wide the choice among them. These bear on the degree of meaning available, which EXI-01-d says varies with circumstance (its Proposition~2.4), and not on its availability.

The minimal form makes the assumption easier to defend than its stronger relatives. It is a much stronger claim that every agent has a \emph{wide} range of worthwhile projects available, or that every agent can attain a \emph{high} degree of meaning, and those claims are false: deprivation narrows the range, and injustice lowers the attainable degree. Assumption~3 does not deny either fact. The comment on \Lm{5.2} returns to the point.

\paragraph{\enquote{Project}, again.} If \enquote{project} were read in the narrow sense of an undertaking with a completion condition, Assumption~3 would be doubtful, as EXI-01-c observes: prisoners, those who labor without respite, the very ill and the very poor may have no such undertaking open to them. Proposition~\ref{p:broad} has argued that the broad sense is the correct reading of Definition~1, and Assumption~3 uses the word in the same sense. In the broad sense, the care of a fellow prisoner, the attention a laborer gives her children and the bedridden patient's attention to those who visit her are projects, provided they are sustained and caring, and they have objective worth because their objects do. EXI-01-c derives Assumption~3, so read, from its Assumption~2.2 and its Lemma~2.3, and so makes it a theorem rather than a premise.

\lem{4.4}{in which she is able to engage}

\paragraph{The modality.} \enquote{Able} is a modal word, and several modalities are candidates. A project is \emph{physically} open to an agent if nothing in the world prevents her from undertaking it. It is \emph{epistemically} open if she knows of it or could come to know of it. It is \emph{motivationally} open if she could come to care about it. It is \emph{socially} open if she is permitted and enabled by those around her to undertake it. The text does not say which is meant.

The best reading is the conjunction, restricted to what the agent can do from where she is. To be able to engage in a project, in the sense that matters for availability, is to be able, from one's actual situation, to come to act with care in ways directed at it. That requires that the project be physically within reach, that the agent be able to come to know of it, and that she be able to come to care about it, though not that she already knows of it or cares about it. The last condition is the one most often overlooked. A person may be surrounded by worthwhile objects and unable to care about any of them, as in severe depression; and Assumption~3 then requires that the inability be temporary or partial. The comment on \Lm{2.5} argued that the depressed agent who acts with care but without feeling is engaged; the agent who cannot act with care at all is a harder case, and Volume~III, Part~A, discusses it. Here we note only that \enquote{able} is a dispositional notion, and that a disposition can be masked without being lost.

\paragraph{\enquote{she}.} The pronoun is the generic feminine used throughout the discipline. It has no philosophical significance, but it does have a grammatical one: it is bound by \enquote{every agent} and shows that the ability is the agent's own. The availability claim is not that projects of objective worth are available to the agent in the sense that someone else could engage in them on her behalf, but that she can.

\paragraph{\enquote{engage}.} The verb is \enquote{engage}, without \enquote{actively}. The text is not careless here: the ability in question is the ability to engage, and active engagement in the sense of Definition~1 is its exercise. The proof of Theorem~4 must move from the ability to engage to a life that consists in active engagement, and the comment on \Lm{7} examines how it does so.

\paragraph{Frankl's third way.} \citet{Frankl1962} held, from his experience of the camps, that meaning can be found in three ways: by creating a work or doing a deed, by experiencing something or encountering someone, and by the attitude one takes toward unavoidable suffering. The first two correspond to engagement with projects in the broad sense. The third raises a question for Assumption~3. If an agent is so placed that no project, not even the care of another, is open to her, can she still find meaning in the way she bears her condition? On Definition~1, she can only if bearing one's condition with courage or dignity is a project of objective worth in which she is able to engage. We think it is: the exercise of courage is directed activity, it can be sustained and caring, and courage is a good that idealized reflection endorses. If so, Frankl's third way is a limiting case of Definition~1 rather than an exception to it, and it is the case that makes Assumption~3 hardest to refute, since an agent capable of intentional action is always able to take some attitude toward her situation. We record this reading without relying on it: the argument of EXI-01-c establishes Assumption~3 without appeal to it.

\paragraph{Support for the assumption.} The text does not argue for Assumption~3. The family supplies three supports.

The first is EXI-01-c, which derives Assumption~3 from the premise that every agent capable of intentional action can direct some intentional activity, with care, at a worthwhile object of engagement, another person or some part of the world she can come to understand (its Assumption~2.2), together with the lemma that sustained caring activity directed at such an object is a project of objective worth (its Lemma~2.3). The premise is supported by the ubiquity of persons and of the world as an object of understanding. It is weaker than Assumption~3 in what it asserts about projects and stronger in what it asserts about objects, and EXI-01-c argues that it is better supported.

The second is EXI-17-a, which shows that meaning does not require achievement (its Theorem~3.3) and that it is therefore not hostage to luck in outcomes (its Corollary~3.6). This removes one of the main threats to Assumption~3. If meaning required success, the projects in which an agent could engage would confer meaning only if they could succeed, and many agents have no project that can succeed. Since meaning requires only engagement, a project can confer meaning though it fails, and the projects open to an agent are not reduced by the world's refusal to cooperate.

The third is EXI-04-a's Lemma~3.1, that every agent capable of intentional action takes some ends seriously. This does not entail Assumption~3, since the ends one takes seriously may be worthless. But it shows that the subjective capacity on which engagement depends is present in every agent in the domain, so that a failure of Assumption~3 would have to be a failure of worth or of reach, not a failure of the capacity for engagement as such.

\paragraph{Objections.}

\begin{objection}[Radical deprivation]\label{o:depriv}
There are agents capable of intentional action who are so deprived, by imprisonment, poverty, illness or oppression, that nothing worthwhile is within their reach. Assumption~3 is false for them.
\end{objection}

\begin{reply}
The objection is the one EXI-01-c was written to answer, and it is answered by the broad reading of \enquote{project}. The prisoner can care for a fellow prisoner and can come to understand something; the exhausted laborer can attend to her children; the bedridden patient can attend to her visitors and to what she hears. These are projects of objective worth in the sense of Definition~1, and they are within reach. What deprivation takes away is the range of projects and the degree of meaning attainable, which EXI-01-d acknowledges (its Proposition~2.4). It does not take away the possibility of meaning, because the objects of worth that remain in every human situation, persons and the world as an object of understanding, are not things that deprivation can remove.
\end{reply}

\begin{objection}[Triviality]\label{o:trivialA3}
On the broad reading, Assumption~3 is trivially true, and Theorem~4 inherits its triviality. The text purchases universality by emptying \enquote{project} of content.
\end{objection}

\begin{reply}
The assumption is not trivial on any reading. It requires, for each agent, an object of objective worth, which a subjectivist denies exists; it requires that the agent be able to engage with it, which requires that it be within her physical, epistemic and motivational reach; and it requires, through the word \enquote{project}, that the engagement be sustained and caring, not momentary. What the broad reading removes is the requirement that a project have a completion condition, a public product or an institutional form, and these are requirements that, as EXI-17-a shows for achievement, have nothing to do with meaning. A claim that is easy to satisfy because its conditions are the right ones is not thereby trivial.
\end{reply}

\subsection{Hard cases}\label{s:hard}

The objection from radical deprivation is usually pressed with examples, and an assumption with a universal quantifier should be tested on the hardest of them. We take seven. In each we ask whether the agent, while she remains capable of intentional action, has a project of objective worth within her reach, and we read Assumption~3 as indexed to times, as the comment on \Lm{5.3} argues. A fuller treatment, with the empirical literature on each condition, belongs to Volume~III, Part~A; the purpose here is to show how the lemmata apply.

\paragraph{Solitary confinement.} The prisoner in solitary confinement lacks the first of EXI-01-c's two routes, care for another person present to her, for long stretches. Engagement with absent persons through memory and imagination is real attention, but Principle~P2 of EXI-17-a requires that engagement bear on its object, and it is doubtful that thinking of an absent friend bears on the friend. The second route remains: some part of the world can be attended to and understood. Prisoners in solitary confinement have worked through mathematics in their heads, reconstructed and studied what they remembered, learned the habits of insects in their cells; what is attended to is real, and understanding of it is a good that does not depend on the prisoner's attitudes. The project is meager, and the degree of meaning it can sustain is low. But it is a project of objective worth in the broad sense of Proposition~\ref{p:broad}, and the prisoner is able to engage in it. The case confirms the importance of the second route, which EXI-01-c includes precisely because the first can be taken away.

\paragraph{Severe chronic pain.} Pain consumes attention, and engagement requires attention. A person in severe chronic pain may be unable, for long periods, to sustain any directed activity. Two features of the lemmata bear on the case. First, \enquote{able} in \Lm{4.4} is dispositional, and the disposition to engage is masked, not lost, by pain: in the intervals in which the pain recedes, the person can engage, and the projects she takes up then are hers. Second, Proposition~\ref{p:neutral} requires only that engagement be sustained over a non-null part of the life, not over every part. The availability claim holds; the degree attainable is reduced, which is EXI-01-d's Proposition~2.4.

\paragraph{Advanced dementia.} Dementia erodes the capacity for intentional action gradually. In its earlier and middle stages, the person retains the capacity to act for reasons and to care about others, and relationships of affection with those who care for her remain within her reach and have worth. Assumption~3 holds. In its final stages, the capacity for intentional action may be lost, and the person then falls outside the domain of \Lm{4.2}. The availability claim does not say that meaning is available to her then, and it does not say that it is not; the comment on \Lm{4.2} has noted that the question whether lives outside the domain can have meaning in another sense is left open. What the indexed reading secures is that the boundary of availability is the boundary of agency, not a boundary set by the illness before agency is lost.

\paragraph{Early childhood.} A young child acquires the capacity for intentional action gradually, and with it the capacity for engagement. Play and the exploration of the world are the child's first engagements, and much of them is engagement with understanding and with other persons, which have worth. The domain of Assumption~3 grows with the child, and within it the assumption holds. Whether play as such is a project of objective worth is a further question, which the case does not require us to settle, since the child's play typically includes engagement with persons and with understanding.

\paragraph{Forced labor and enslavement.} The person whose labor is owned by another and directed to another's ends has the narrowest range of projects that human societies have imposed. The ends to which her labor is directed are not hers, and whether they have worth is irrelevant to the meaning of her life, since she is not engaged with them; EXI-11-a's principle that worth is not conferred by assignment (its Principle~P1) has an analogue here, that meaning is not conferred by assignment either. But the persons among whom she lives remain, and with them the care of children and kin, friendship, the transmission of what one knows, and the pursuit of fairness in the form of resistance. These are projects of objective worth, and they have been pursued under the worst conditions that history records. The case is the strongest illustration of EXI-01-d's insistence that availability does not imply indifference: that meaning remains available under oppression is no reason to tolerate oppression, which reduces the degree of meaning as it reduces every other good.

\paragraph{The last days.} A person who is dying, with days to live and little strength, has a short continuation before her. The indexed reading of Theorem~4 requires that some continuation available to her contain engagement with a project of objective worth to a positive extent. The last conversations of the dying with those they love are such engagement: directed, caring and bearing on real persons. On the proportional reading of the measure, they make up a positive share of the remaining life; on the absolute reading, they add a positive amount. The case shows why the indexed reading matters: it is not too late for the dying to engage, and Definition~1 counts what they do.

\paragraph{Severe depression.} The hardest case is the one in which the capacity to care is itself impaired. The depressed nurse of \Lm{2.5} acts with care though without feeling; a person in a severe depressive episode may be unable to act with care at all, not because nothing worthwhile is within reach but because nothing moves her. Assumption~3 then holds only if \enquote{able to engage} is read dispositionally, as the comment on \Lm{4.4} argued: the capacity to care is masked by the illness, and typically returns with treatment or with time. If there were a person in whom the capacity to care was permanently and wholly lost while the capacity for intentional action remained, she would be a counterexample to Assumption~3 as read here. We doubt that such a case is possible, since acting for reasons involves caring about something, but we record it as the point at which the assumption is most exposed, and Volume~III, Part~A, examines it.

\begin{proposition}[Reach in the hard cases]\label{p:reach}
In each of the seven cases, while the agent remains capable of intentional action and her capacity to care is not permanently lost, some project of objective worth is within her reach, and Assumption~3 holds of her at that time.
\end{proposition}

\begin{proof}[Argument]
By the discussion of each case: the second route of EXI-01-c (understanding some part of the world) remains in solitary confinement; the dispositional reading of \enquote{able} covers chronic pain and depression; the first route (care for persons present) remains in dementia short of the loss of agency, in childhood once agency begins, under forced labor and in the last days.
\end{proof}

The proposition does not say that the meaning available in these cases is great. In most of them it is small, and the circumstances that make it small are evils. It says that it is not zero, and that is what Theorem~4 claims.
